\section*{Project interfaces}
The project exposes the entries below. Each pointer is a full MIPS64 address. The input
stack is eight-byte aligned. Every entry preserves all bits of \texttt{s0} through
\texttt{s7}, \texttt{fp} and \texttt{sp}, returns through the supplied \texttt{ra},
and produces no console output. Submission-owned read-only constants and private stack
storage are allowed. Caller objects named below are disjoint.

\begin{longtable}{p{0.32\linewidth}p{0.61\linewidth}}
\caption{Function inputs and writable caller storage.}\\
\toprule
Entry & Arguments and effect\\
\midrule
\endfirsthead
\toprule Entry & Arguments and effect\\\midrule
\endhead
\texttt{sha256\_schedule} & \texttt{a0}: 64 source bytes; \texttt{a1}: 256-byte,
word-aligned schedule. Write its 64 words. Preserve the source. Return zero in \texttt{v0}.\\
\texttt{sha256\_compress} & \texttt{a0}: eight word-aligned state words;
\texttt{a1}: 64 source bytes; \texttt{a2}: 256-byte, word-aligned workspace.
Replace the state with its compressed result. Only state and workspace are writable.
Return zero in \texttt{v0}.\\
\texttt{sha256\_init} & \texttt{a0}: an eight-byte-aligned 368-byte context.
Initialize its state, byte count and used count independently of old contents.
Return zero in \texttt{v0}.\\
\texttt{sha256\_update} & \texttt{a0}: initialized context; \texttt{a1}: source;
\texttt{a2}: byte count, zero through 4096. Consume those bytes, writing only the
context. Zero count reads no source bytes. Return zero in \texttt{v0}.\\
\texttt{sha256\_final} & \texttt{a0}: context; \texttt{a1}: 32 output bytes.
Finalize and write the digest in big-endian byte order. Only context and output are
writable. Return zero in \texttt{v0}.\\
\texttt{sha256\_file} & \texttt{a0}: terminated path; \texttt{a1}: scratch bytes;
\texttt{a2}: capacity, one through 256; \texttt{a3}: 32 output bytes. Open read-only,
hash through successful EOF and close. Return status zero, one for open failure, or
two for read failure in \texttt{v0}. Preserve output on failure.\\
\bottomrule
\end{longtable}

Source blocks, update inputs, output bytes and scratch buffers may be unaligned.
Access only the selected intervals of caller storage. A finalizer may change any bytes
inside its context, but it must preserve other caller objects. A file routine uses a
private context and may change only the selected scratch interval and, on success,
the 32-byte output interval. Close a descriptor after a read error; an unsuccessful
open creates no descriptor to close.

An update length is limited to 4096 bytes per call for this teaching interface. Several
calls may form a longer message. The accumulated byte count must remain below
\(2^{61}\), so its bit count fits in the standard's 64-bit length field. Inputs outside
the stated domain are not part of this contract. Finalization consumes the context:
initialize it again before starting another message. Repeated finalization is not
specified. Two distinct initialized contexts must remain independent when calls are
interleaved.
